\section{Introduction}
\label{sec:intro}

A transformer trained on sequences drawn from a distribution whose latent structure is resampled
for every sequence cannot memorise that structure in its weights; it must infer it from the
context. The circuit that performs this inference is the \emph{induction head}
\citep{elhage2021,olsson2022}, and a substantial body of work now characterises how it arises
\citep{bietti2023,edelman2024,nichani2024}. Across this literature the informative context is
adjacent, or nearly so: the head matches on the immediately preceding token, or selects among a
small set of candidate lags.

Real sequences are not so tidy. Between two causally linked symbols there may be an arbitrary
amount of material that carries no information about the transition at all. A model that can only
match at a fixed lag is blind to structure separated by such material, and no amount of training
data will teach it otherwise if the mechanism cannot express the dependency.

We study the simplest setting in which this is forced. Sequences are generated by a Markov chain
whose transition table is resampled for every sequence, so nothing about it can be stored in the
weights. With probability $1-\alpha$ the emitted symbol is drawn from a disjoint
\emph{distractor} sub-vocabulary and the chain does not advance. Predicting the next symbol
therefore requires \emph{state recovery}: skipping a run of distractors whose length is
geometric with mean $(1-\alpha)^{-1}$ in order to locate the last state-bearing symbol, and only
then performing the induction match. A single parameter, $\alpha$, interpolates between a task
dominated by distraction and the noiseless setting studied previously.

Our central finding is that gradient descent builds this circuit only when distractors are
\emph{common}. Sweeping $\alpha$, the rate at which a trained model predicts a non-successor
immediately after a distractor stays below $1.2\%$ for $\alpha\le0.7$ and jumps to $51.6\%$ at
$\alpha=0.8$: a factor of $143$ across one step of the sweep, in the direction of \emph{less}
noise.\footnote{The direction is worth emphasising. At $\alpha=0.8$ the task is easier by every
aggregate measure --- the cross-entropy is lower and the Bayes floor is lower --- yet the model is
qualitatively worse at the one sub-task that distinguishes this setting from the noiseless one.}
When distraction is rare the model does not learn to handle it; it settles for copying the
previous symbol, which is correct a fraction $\alpha$ of the time, and abandons state recovery
entirely.

To explain this we analyse the learning dynamics from zero initialisation, where the model
computes the in-context unigram predictor and every quantity is available in closed form. Three
facts emerge. The first layer has \emph{identically zero} gradient, so the circuit cannot form in
both layers at once. The second layer's gradient is structured, but only in the block that reads
the recovered state: because the transition table is resampled, the marginal bigram statistic is
uninformative and the corresponding block receives no diagonal signal in expectation. In-weights
learning is therefore not merely suboptimal here but \emph{gradient-inaccessible}, and the first
descent direction available to the model is already an in-context one. Finally, the second layer
recruits the first, which is then driven towards state recovery.

\paragraph{Contributions.}
\begin{itemize}
\item We introduce in-context learning of Markov structure \emph{under distraction}, a
one-parameter family that contains the standard noiseless setting as its $\alpha\to1$ boundary,
and give a two-layer construction solving it (Section~\ref{sec:construction}).
\item We report a sharp phase transition in distractor density: state recovery is learned for
$\alpha\le0.7$ and not learned for $\alpha\ge0.8$ (Section~\ref{sec:transition}).
\item We analyse the dynamics from zero initialisation and show that the in-weights channel
receives no gradient, that the layers must form sequentially, and that the second layer recruits
the first (Section~\ref{sec:dynamics}, proofs in Appendix~\ref{app:dynamics}).
\item We show our construction and analysis reduce to the noiseless case as $\alpha\to1$,
recovering prior results as a boundary case (Section~\ref{sec:reduction}).
\end{itemize}

\section{Setup}
\label{sec:setup}

\paragraph{Task.}
Fix a vocabulary $\mathcal{V}=\mathcal{S}\cup\mathcal{D}$ partitioned into $V_s$ \emph{signal}
symbols and $V_d$ \emph{distractors}. For each sequence we draw a transition table $\pi$
independently, assigning to every $s\in\mathcal{S}$ a successor set
$\mathrm{succ}_\pi(s)\subset\mathcal{S}$ of size two, and generate $v_{0:T-1}$ by
\begin{equation}
v_t=\begin{cases}
\text{uniform from }\mathrm{succ}_\pi(\mathrm{st}(t\!-\!1)), & \text{w.p. }\alpha
\quad\text{(state advances)},\\[2pt]
\text{uniform from }\mathcal{D}, & \text{w.p. }1-\alpha
\quad\text{(state unchanged)},
\end{cases}
\end{equation}
where $\mathrm{st}(t)$ denotes the chain state after position $t$. The target is
$y_t=v_{t+1}$ and the loss is averaged over the second half of the sequence,
$t\in\mathcal{T}=[T/2,T)$, so that every scored position has a context long enough to identify
the table. Because $\pi$ is resampled, no function of the weights alone can predict better than
the marginal; all of the achievable performance is in-context.

\paragraph{The state and why it is not the previous symbol.}
Distractors do not advance the chain, so the state at position $t$ is the most recent
\emph{signal} symbol at or before $t$. The number of distractors separating a symbol from its
causal parent is geometric with mean $(1-\alpha)^{-1}$ and is unbounded. A previous-token head,
or a selector over any finite set of lags, cannot express this dependency.

\paragraph{Architecture.}
The input embedding is a concatenation of two one-hots and is not learned,
$x_t=[\,e_{v_t}\,;\,e_t\,]\in\mathbb{R}^{d}$ with $d=V+T$. Two attention layers concatenate their
output onto the residual stream,\footnote{Concatenation rather than addition keeps the two layers'
contributions in disjoint coordinates, which is what makes the block decomposition of
Section~\ref{sec:dynamics} exact rather than approximate.} and both use a strictly causal mask.
With combined query--key matrices $W^{(1)}\in\mathbb{R}^{d\times d}$,
$W^{(2)}\in\mathbb{R}^{2d\times 2d}$,
\begin{equation}
h_t=\sum_{u<t} a^{(1)}_{t,u}\,e_{v_u},\qquad
x^{(2)}_t=[\,x_t\,;\,h_t\,],\qquad
f_t=\sum_{u<t} a^{(2)}_{t,u}\,e_{v_u},
\label{eq:arch}
\end{equation}
where $a^{(i)}_{t,\cdot}=\mathrm{softmax}\bigl(x^{(i)}_tW^{(i)}x^{(i)\top}_{\cdot}/\sqrt d+M\bigr)$
and the value circuits are fixed to the selector that passes the symbol one-hot. Appendix
\ref{app:gauge} shows that fixing the first layer's value circuit costs no generality --- it is a
choice of coordinates --- while fixing the second layer's is a genuine restriction, which we
state as such.

\paragraph{Readout.}
Since $f_t$ in \eqref{eq:arch} is a convex combination of one-hots it is already a distribution,
with $\sum_w f_t[w]=1$ identically. We therefore read out by renormalisation against a fixed
bias, with no unembedding, no logarithm and no output softmax:
\begin{equation}
p_t=\frac{f_t+b}{\sum_w\bigl(f_t[w]+b_w\bigr)},\qquad
b_w=\begin{cases} b_s, & w\in\mathcal{S},\\ b_d, & w\in\mathcal{D},\end{cases}
\qquad
b_d=\frac{(1-\alpha)(1+V_s b_s)}{V_d\,\alpha}.
\label{eq:readout}
\end{equation}
The choice of $b_d$ is forced: it is the unique value for which a model with a perfect distractor
filter assigns exactly $(1-\alpha)/V_d$ to each distractor. Appendix~\ref{app:readout} derives it
and shows that this readout is what an unembedding followed by a softmax computes when the
intervening nonlinearity is $\log(f+b_s)$.

\begin{table}[!ht]
\caption{Notation.}
\label{tab:notation}
\begin{center}
\begin{tabular}{ll}
\toprule
$\mathcal{S},\mathcal{D}$ & signal and distractor sub-vocabularies, sizes $V_s$, $V_d$ \\
$\alpha$ & probability that a position emits a signal symbol and advances the chain \\
$\pi$ & per-sequence transition table; $\mathrm{succ}_\pi(s)$ its successor set \\
$\mathrm{st}(t)$ & chain state after position $t$; the most recent signal symbol \\
$h_t, f_t$ & output of layer 1 and layer 2 at position $t$, both in the simplex \\
$m_t$ & running symbol frequency $\tfrac1t\sum_{u<t}e_{v_u}$ of the context \\
$P,Q,R,C$ & the four $d\times d$ blocks of $W^{(2)}$ (Section~\ref{sec:dynamics}) \\
\bottomrule
\end{tabular}
\end{center}
\end{table}

\section{The state-recovering induction head}
\label{sec:construction}

We first establish that the task is solvable by this architecture, and how. The construction
divides the work between the layers.

\begin{theorem}[Construction]
\label{thm:construction}
There exist weights $W^{(1)},W^{(2)}$ such that, up to softmax saturation, layer 1 places all of
its attention on the most recent signal position and layer 2 outputs the empirical successor
frequency of the current state:
\begin{equation}
h_t=e_{\mathrm{st}(t)},\qquad
f_t[w]=\frac{\#\{u<t:\ \mathrm{st}(u\!-\!1)=\mathrm{st}(t),\ v_u=w\}}
{\#\{u<t:\ \mathrm{st}(u\!-\!1)=\mathrm{st}(t)\}} .
\end{equation}
Composed with the readout \eqref{eq:readout}, the model realises add-$b_s$ smoothed
context-Bayes prediction.
\end{theorem}

The mechanism, proved in Appendix~\ref{app:construction}, uses three devices. First, layer 1
assigns score $-\infty$ to every key whose symbol is a distractor, so the attention is supported
on signal positions only. Second, among those it applies a recency ramp linear in $t-u$, so the
softmax concentrates on the most recent one; the ramp is what makes the unbounded lag expressible
with a fixed number of parameters. Third, layer 2 compares the query's own symbol against the key's
recovered state $h_u$, and averages the symbols it finds --- turning an attention average into a
count, since the attended set is exactly the set of past transitions out of the current state.

Two remarks. The first layer is \emph{not} a previous-token head: it is a previous-\emph{signal}-token
head, and the distinction is the whole content of the construction. And the query in layer 2
matches on its own symbol when that symbol is a signal, falling back to the recovered state when it
is a distractor; a single head suffices because those two branches ask the same question.

\section{A distractor-density phase transition}
\label{sec:transition}

Theorem~\ref{thm:construction} says the circuit exists. It does not say gradient descent finds it,
and the central empirical result of this paper is that whether it does depends sharply on
$\alpha$.

We train the model of Section~\ref{sec:setup} at each $\alpha$ and then measure, on held-out
sequences and using each sequence's own table $\pi$, the rate at which the model predicts a signal
symbol that is \emph{not} a successor of the current state. We report this separately at
post-distractor positions --- where the state must be recovered by skipping backwards --- and at
post-signal positions, where the state is directly visible.

\begin{table}[!ht]
\caption{Post-distractor error: the rate at which the trained model predicts a signal symbol that
is not a successor of the current state, at positions immediately following a distractor.
Held-out sequences, each scored against its own transition table. The change between
$\alpha=0.7$ and $\alpha=0.8$ is a factor of $143$.}
\label{tab:transition}
\begin{center}
\begin{tabular}{lccccccccc}
\toprule
$\alpha$ & 0.1 & 0.2 & 0.3 & 0.4 & 0.5 & 0.6 & 0.7 & 0.8 & 0.9 \\
\midrule
post-distractor error & 0.89\% & 0.31\% & 0.27\% & 1.16\% & 0.51\% & 0.44\% & 0.36\%
  & \textbf{51.59\%} & \textbf{49.94\%} \\
\bottomrule
\end{tabular}
\end{center}
\end{table}

The transition is abrupt and has no intermediate value in our sweep. Below it the model is
essentially never wrong about which symbols may follow the current state; above it, it is wrong
half the time --- close to what a model that ignores the state entirely and reports the marginal
would achieve.

\paragraph{Interpretation.}
The aggregate cross-entropy does not reveal this. It decreases monotonically in $\alpha$, because
distractors are individually easy to predict (they are uniform on $\mathcal{D}$, and the readout
bias $b_d$ supplies exactly the right mass). The failure is confined to the sub-task that the
aggregate barely weights. A model that has not learned state recovery can still do well overall by
copying the previous symbol, which is the correct state a fraction $\alpha$ of the time; at
$\alpha=0.8$ that shortcut is right four times in five, and the remaining fifth is too small a
share of the loss to force the harder circuit into existence.

% TODO: add the post-signal control column. The claim "state recovery specifically fails" requires
% the post-signal error to stay flat while the post-distractor column explodes. Our current
% measurement of the control had an off-by-one in the state update and is not reportable.
% TODO: seed bands. Table 1 is a single seed.
% TODO: confirm the transition replicates in the reduced-vocabulary configuration used for
% Section 5; Table 1 is measured in the larger configuration (Appendix D).

\section{Learning dynamics from zero initialisation}
\label{sec:dynamics}

To explain Table~\ref{tab:transition} we ask what gradient descent does at the very beginning,
where everything is computable. Initialise $W^{(1)}=W^{(2)}=0$. Both softmaxes are then uniform
over the admissible keys and, by \eqref{eq:arch},
\begin{equation}
h_t=f_t=m_t:=\frac1t\sum_{u<t}e_{v_u},
\end{equation}
the running symbol frequency: the model begins as the in-context unigram predictor.

Split $W^{(2)}$ into $d\times d$ blocks according to $x^{(2)}=[x;h]$, so that the score
decomposes into four interpretable terms,
\begin{equation}
W^{(2)}=\begin{bmatrix}P & Q\\ R & C\end{bmatrix},\qquad
\sqrt d\,S^{(2)}_{t,u}=\underbrace{x_tPx_u^\top}_{\text{symbol--symbol}}
+\underbrace{x_tQh_u^\top}_{\text{induction}}
+\underbrace{h_tRx_u^\top}_{\text{cross}}
+\underbrace{h_tCh_u^\top}_{\text{state--state}}.
\label{eq:blocks}
\end{equation}
The symbol--symbol block $P$ is the in-weights channel: it scores a key by its own symbol against
the query's, which is what a model would use to store a fixed bigram table. The induction block
$Q$ is the in-context channel: it scores a key by the \emph{state} that preceded it.

\begin{theorem}[Dynamics at initialisation]
\label{thm:dynamics}
At $W^{(1)}=W^{(2)}=0$, writing $\omega_t=(m_t[y_t]+b_{y_t})^{-1}>0$:
\begin{enumerate}
\item[(i)] \textup{(Gating.)} $\partial L/\partial W^{(1)}=0$ identically. The first layer cannot
move until the second has left the origin.
\item[(ii)] \textup{(Closed form.)}
$\displaystyle -\frac{\partial L}{\partial W^{(2)}}
=\frac{1}{|\mathcal{T}|\sqrt d}\sum_{t\in\mathcal{T}}\frac{\omega_t}{t}\sum_{u<t}
\bigl(\mathbf 1[v_u=y_t]-m_t[y_t]\bigr)\,x^{(2)}_t\otimes x^{(2)}_u .$
\item[(iii)] \textup{(The in-weights channel is dead.)} In expectation over the resampled table,
the induced update to $P$ has no diagonal component, because
$\mathbb{E}_\pi[\mathbf 1[w=y_t]\mid \mathrm{st}(t)=v]=\alpha/V_s$ independently of
$w\in\mathcal{S}$. The corresponding update to $Q$ does have one.
\end{enumerate}
\end{theorem}

Part (iii) is the reason the title of this section is about in-context learning. The marginal
bigram statistic is uninformative \emph{by construction of the task}: averaged over the resampled
table, every signal symbol is equally likely to follow every other. The model cannot take a step
towards storing a transition table even if such a step would help, because the gradient in that
direction is zero. The first direction available to it is the one that reads the context.

Measured at $\alpha=0.6$, the diagonal-minus-off-diagonal statistic of the first descent
direction, restricted to the signal sub-block, is $-1.48\times10^{-4}$ on $P$ against
$+9.62\times10^{-5}$ on $Q$ and $+9.80\times10^{-6}$ on $C$; the cosine between the induction
block and the construction's is $+0.867$, against a random-matrix baseline of $0.010$. The first
gradient does not merely avoid the wrong direction --- it points most of the way to the right one.

\paragraph{Steps one and two.}
Appendix~\ref{app:dynamics} continues the analysis. After one step the first layer's gradient
becomes non-zero and is proportional to those blocks of the second layer's update that contract
with $h$, so the second layer \emph{recruits} the first
(Lemma~\ref{lem:step1}). After two steps the first layer is driven towards suppressing distractor
keys and sharpening onto the most recent signal position, which is state recovery
(Lemma~\ref{lem:step2}). The circuit is built outwards from the readout rather than inwards from
the input.

\paragraph{Connecting to the transition.}
The coefficient in Theorem~\ref{thm:dynamics}(ii) is supported on positions whose target is a
successor of the current state. The positions at which state recovery is \emph{required} --- those
following a distractor --- form a fraction $1-\alpha$ of the scored window. As $\alpha\to1$ that
share vanishes, and with it the gradient pressure to build the first layer at all.
% TODO: make this quantitative. This is the step that turns Section 5 into an explanation of
% Section 4 rather than a plausible story alongside it.

\section{Reduction to the noiseless case}
\label{sec:reduction}

Our setting contains the standard one as a boundary.

\begin{theorem}[Reduction]
\label{thm:reduction}
As $\alpha\to1$ the distractor sub-vocabulary is never emitted; the state at $t$ is $v_t$ itself;
the construction of Theorem~\ref{thm:construction} degenerates to a previous-token head composed
with an induction head; and the task becomes in-context learning of a Markov chain with a
per-sequence transition table, as studied by \citet{edelman2024}.
\end{theorem}

This positions the prior work as our $\alpha=1$ limit rather than as a competing account, and it
provides a consistency check: our measurements at $\alpha=1.0$ should reproduce the qualitative
behaviour reported there, and they do.
% TODO: state which quantities are compared and give the numbers.

\section{Related work}
\label{sec:related}

\paragraph{Induction heads.}
The circuit was identified by \citet{elhage2021} and connected to in-context learning by
\citet{olsson2022}. \citet{bietti2023} analyse a two-layer model in which a global and an
in-context component coexist, which is structurally analogous to our signal/distractor split,
though their distinction is between what is stored and what is inferred rather than between
informative and uninformative positions.

\paragraph{In-context learning of Markov structure.}
\citet{edelman2024} train two-layer attention-only transformers on sequences whose transition
matrix is resampled per sequence and document a multi-phase trajectory: uniform, then in-context
unigram, then a rapid transition to the in-context bigram solution. Our model at zero
initialisation \emph{is} their unigram phase, which we make precise in
Theorem~\ref{thm:dynamics}. Their setting is our $\alpha\to1$ limit
(Theorem~\ref{thm:reduction}); the distractor sub-vocabulary and its density are what is new here.

\paragraph{Training dynamics.}
\citet{nichani2024} prove that gradient descent recovers the causal structure of a latent graph,
and \citet{unveiling2024} prove a three-stage convergence for a two-layer transformer on $n$-gram
chains in which a feed-forward selector chooses among parents at fixed offsets. Both analyse
architectures in which the informative position is at a bounded, enumerable displacement. Our
first layer must instead locate a parent at an unbounded, geometrically distributed lag, and our
contribution is the dependence of that process on the distractor density.

\paragraph{Selecting causal structure.}
\citet{selective2025} study interleaved chains in which the \emph{lag} varies while the
transition probabilities are held fixed, and construct a three-layer selective induction head that
identifies the correct lag in context. The complementary axis is instructive: they vary the lag
with fixed structure, we resample the structure with a lag that is not selected from a finite set
but must be found by skipping uninformative material. A selector over a finite lag set cannot
express a geometric lag distribution.

\section{Discussion and limitations}
\label{sec:limitations}

The phenomenon we report is a failure of a rare sub-task to attract gradient, and the mechanism is
not specific to distractors: whenever a capability is required on a vanishing fraction of the
training distribution, and a shortcut exists that is correct on the remainder, the shortcut is
what training will produce. What makes the present setting useful is that the rare sub-task, the
shortcut and the correct circuit are all explicitly identifiable.

Several limitations should be stated plainly. The task is synthetic, the vocabulary small and the
model two layers deep; we make no claim about trained language models. The value circuits are
held fixed, which Appendix~\ref{app:gauge} shows is free in the first layer and a real restriction
in the second. Theorem~\ref{thm:dynamics}(iii) is proved for the expectation of the in-weights
block; the corresponding lower bound for the in-context block is verified numerically but its
proof requires control of the chain's mixing, which we do not yet have. Finally, weight-space
agreement with the construction is not monotone during training: it peaks within the first ten
steps and then declines while the loss continues to improve. We report this but cannot explain it.
